\documentclass[11pt]{article}

\usepackage{amsmath,amssymb}
\usepackage{booktabs}
\usepackage{graphicx}
\usepackage{hyperref}
\usepackage[margin=1in]{geometry}
\usepackage{natbib}
\usepackage{lineno}
\usepackage{xcolor}

\title{Direction Instability Identifies Drug Mechanism Transport\\
While Geometric Alternatives Measure Magnitude}

\author{Elliot Tower\\
\texttt{elliot@elliottower.ai}}

\date{July 2026}

\linenumbers

\begin{document}
\maketitle

\begin{abstract}
High-throughput perturbation screens produce multi-dimensional signatures
for thousands of drugs and gene knockdowns across panels of cell lines.
Geometric metrics---direction instability, Grassmannian geodesic
distance, profile variability, optimal transport distances---are used to
quantify whether a perturbation acts consistently across cellular
contexts. We show that three of these four metric families are dominated
by estimation quality: any variable that improves the precision of the
estimated object (effect magnitude, sample size) mechanically shifts the
distance metric. Across three perturbation modalities (LINCS L1000,
Perturb-seq, DepMap CRISPR screens) and two measurement technologies
(transcriptomic and morphological), the same artifact produces effects ranging from large (DepMap Cohen's
$d = 2.65$, Perturb-seq geodesic $d = 0.55$) to ambiguous
(cross-cell-type transport $d = -0.14$ to $-0.172$), all of which
invert or vanish after magnitude control. The fourth metric family,
cosine-based direction instability, is immune by construction: it
normalizes by magnitude before comparison. This immunity is functional:
on a pre-registered construction-neutral benchmark, direction instability
predicts cross-cell-line mechanism transport (leave-one-out AUROC $=
0.945$, 95\% CI $[0.937, 0.952]$, $n = 2{,}694$ LINCS drugs),
outperforming all alternatives including Fr\'{e}chet variance ($0.857$,
$[0.841, 0.872]$). On the cosine-matched outcome, Fr\'{e}chet variance
is marginally higher ($0.976$ vs.\ $0.969$, $\Delta = -0.007$, $[-0.012,
-0.002]$)---consistent with the construction-matching critique this paper
advances. Direction instability stratifies drugs by mechanism of action
(machinery vs.\ receptor $d = 1.94$ in LINCS, $d = 0.77$ in JUMP Cell
Painting) and reveals an HDAC inhibitor selectivity gradient visible
only through directional analysis. Three pre-registered experiments
(Perturb-seq cosine transport, construction-neutral LINCS benchmark,
cross-modality MOA stratification) yielded five of ten criteria passing
and five failing, with each failure interpretable. Pre-registered
criteria and analysis code were frozen before experiments (commits
\texttt{1dc20a2}, \texttt{249abaf}, \texttt{8d74bd0}). The corrective
rule: normalize for effect magnitude before comparing groups defined by
effect strength, or geometric ``divergence'' reduces to the tautology
that bigger effects are noisier.
\end{abstract}

\section{Introduction}

High-throughput perturbation screens produce multi-dimensional signatures
for thousands of drugs and gene knockdowns across panels of cell lines.
A natural question follows: does this perturbation act consistently
across contexts, or does it produce divergent effects depending on
cellular background? Geometric metrics formalize this question.
Direction instability measures the mean pairwise angular deviation of
perturbation signatures across cell lines. Grassmannian geodesic
distance measures how much a perturbation's response subspace rotates
between cell types. Profile variability (standard deviation across
contexts) measures raw spread.

These metrics are used to characterize drug mechanism conservation, gene
essentiality, and perturbation quality in LINCS L1000
\citep{subramanian2017next}, Perturb-seq \citep{replogle2022mapping},
JUMP Cell Painting \citep{chandrasekaran2023jump}, and DepMap CRISPR
screens. The implicit assumption is that ``more divergent'' reflects
biology---context-dependent mechanisms, compensatory responses, or
pathway rewiring.

This paper makes two contributions. First, we demonstrate that
variance-based, subspace-based, and distribution-based formulations of
these metrics are dominated by estimation quality---any variable that
improves the precision of the estimated object (effect magnitude, sample
size) mechanically shifts the distance metric. Cosine-based metrics
(direction instability) are immune because they normalize by magnitude
before comparison. Second, we exploit this immunity to show that
direction instability captures genuine biological signal: it predicts
cross-cell-line mechanism transport in LINCS L1000 (leave-one-out AUROC
$= 0.945$, 95\% CI $[0.937, 0.952]$ on a construction-neutral outcome),
stratifies drugs by mechanism of action across two measurement
technologies, and reveals an HDAC inhibitor selectivity gradient that
variance-based metrics cannot isolate.

Our evidence comes from a sequence of pre-registered and exploratory
experiments across three data modalities, reported with an explicit
integrity boundary between confirmatory and post-hoc analyses.
The discovery unifies two previously puzzling observations: (1)
biological confound corrections fail to rerank direction instability
across three domains (Spearman $\rho = 0.83$--$0.99$), and (2) essential
gene knockdowns appear to ``diverge'' more than non-essentials ($d =
0.55$ in Perturb-seq). The first reflects genuine directional stability
of perturbation signatures. The second reflects magnitude-confounding of
the comparison. Both dissolve under the same DAG.

\section{Data}

\subsection{LINCS L1000}

We used Level 5 z-score profiles from LINCS L1000
\citep{subramanian2017next} (GEO accession GSE92742), covering 978
landmark genes across up to 70 cell lines. After filtering to drugs
tested in $\geq 3$ cell lines, the dataset contains 8{,}949 compounds.
Drug annotations (mechanism of action, target, clinical phase) were
obtained from the Drug Repurposing Hub \citep{corsello2017drug}. For the
benchmark analysis (Section~\ref{sec:benchmark}), the threshold was
raised to $\geq 10$ cell lines, yielding 2{,}694 drugs.

\subsection{Perturb-seq}

We used genome-scale Perturb-seq data from \citet{replogle2022mapping},
covering 1{,}676 gene knockdowns in K562 and RPE1 cell types. Pseudobulk
signatures were computed as mean expression across cells sharing the same
guide, restricted to genes with $\geq 20$ cells in both cell types.
Subspace data (top-5 PCA subspaces per gene per cell type) were extracted
from the same dataset for geodesic distance and EMD analyses.

\subsection{DepMap CRISPR screens}

We used Chronos dependency scores from the DepMap 24Q4 release (1{,}178
cell lines). Essential genes were defined by DepMap's common-essentials
list ($n = 1{,}242$); non-essential genes by the nonessentials list
($n = 726$). PC1 was removed before computing variability to control for
global fitness effects (PC1 removal changed the effect by 0.2\%).

\subsection{JUMP Cell Painting}

We used morphological profiles from the JUMP Cell Painting dataset
\citep{chandrasekaran2023jump}, covering 25{,}254 compounds across
3{,}180 features in U2OS cells. Compounds were matched to LINCS drugs by
InChIKey prefix (first 14 characters), yielding 330 matched compounds
with MOA annotations.

\section{Metrics and Their Magnitude Sensitivity}

\subsection{Direction instability (cosine-based, magnitude-invariant)}

For a perturbation tested in $K$ contexts with signature vectors
$\mathbf{s}_1, \ldots, \mathbf{s}_K \in \mathbb{R}^d$:
\begin{equation}
  D = 1 - \frac{2}{K(K-1)} \sum_{i < j}
    \frac{\mathbf{s}_i^\top \mathbf{s}_j}
         {\|\mathbf{s}_i\| \, \|\mathbf{s}_j\|}
\end{equation}
The $\ell_2$ normalization makes $D$ invariant to the magnitude of each
$\mathbf{s}_i$. A drug that produces a strong, consistent signal across
cell lines receives the same $D$ as one producing a weak, consistent
signal.

\subsection{Profile variability (variance-based, magnitude-sensitive)}

For a gene with dependency scores $x_1, \ldots, x_K$ across $K$ cell
lines:
\begin{equation}
  V = \text{SD}(x_1, \ldots, x_K)
\end{equation}
This metric grows with the absolute magnitude of effects. A gene with
scores uniformly near $-1$ (strongly essential everywhere) and one with
scores clustered near $0$ (non-essential everywhere) will have different
$V$ even if both are equally ``consistent''---the essential gene's
scores simply occupy a larger range.

\subsection{Grassmannian geodesic distance (subspace-based, indirectly magnitude-sensitive)}

For a gene knocked down in two cell types, let $U_1, U_2 \in
\text{Gr}(k, d)$ denote the top-$k$ PCA subspaces of the perturbation
response. The geodesic distance $\delta(U_1, U_2) =
\|\boldsymbol{\theta}\|_2$ (principal angles) is geometrically
magnitude-invariant. However, $U_1$ and $U_2$ are \emph{estimated} from
finite samples. Perturbations with larger effect sizes produce
higher-SNR data, yielding better-estimated subspaces. This estimation
quality difference propagates into the distance metric, creating an
indirect magnitude sensitivity that is absent from the population-level
definition but present in every finite-sample application.

\subsection{Complementary metrics}

In addition to direction instability, we computed three complementary
metrics for each LINCS drug: magnitude coefficient of variation
($\text{CV} = \text{SD}(\|\mathbf{s}_i\|) / \text{mean}(\|\mathbf{s}_i\|)$),
Fr\'{e}chet variance (trace of the signature covariance matrix), and
mean top-gene Jaccard similarity (overlap of the top 50 most perturbed
genes across cell-line pairs). These serve as comparison baselines for
the multi-metric benchmark (Section~\ref{sec:benchmark}).


\section{Part I: The Magnitude Confound}

\subsection{Pre-registration integrity chain}

All pre-registrations for the magnitude confound experiments were
committed before the author reviewed experimental output (commit
\texttt{249abaf}). The DepMap replication hypothesis was committed
separately (commit \texttt{98897ed}) before execution. The analysis
script was frozen (commit \texttt{92e3276}) before running. The
magnitude-confound analysis (Section~\ref{sec:artifact}) is exploratory
and labeled as such.

\subsection{Experiment: biological corrections do not rerank direction instability}

\textbf{Hypothesis:} Removing confounding features (toxicity genes in
LINCS, essential-gene subspace in Perturb-seq, cell-health features in
JUMP-CP) would substantially reorganize the direction instability ranking
($\rho < 0.70$).

\textbf{Decision criteria:} $\rho < 0.70$: correction bites.
$0.70 \leq \rho < 0.83$: marginal. $\rho \geq 0.83$: does not bite.

\textbf{Result:} Corrections do not bite in any domain (Table~\ref{tab:robustness}).

\begin{table}[h]
\centering
\caption{Rank preservation after biological corrections across three
perturbation domains. All Spearman correlations exceed the pre-registered
``doesn't bite'' threshold of $\rho \geq 0.83$.}
\label{tab:robustness}
\begin{tabular}{llcc}
\toprule
Domain & Correction & $\rho$ & $N$ \\
\midrule
LINCS (global) & Toxicity genes removed & $>0.99$ & 8{,}949 \\
LINCS (HDAC gradient) & Toxicity genes removed & $0.83$ & 20 \\
Perturb-seq & Essential-gene subspace removed & $0.91$ & 1{,}676 \\
JUMP-CP & Cell-health features removed & $0.99$ & 25{,}254 \\
\midrule
Same-tissue LINCS & Within-tissue variance genes & $0.98$--$0.99$ & 176--358 \\
\bottomrule
\end{tabular}
\end{table}

This robustness is independently valid because direction instability is
cosine-based and magnitude-invariant. The ranking is robust because the
underlying geometry is stable, which we now understand is partially
because cosine normalization removes the primary source of apparent
variation (magnitude).

\subsection{Experiment: essential-gene divergence replicates in DepMap}

\textbf{Exploratory prior (disclosed):} In Perturb-seq, essential gene
knockdowns showed higher Grassmannian geodesic distance than non-essentials
(Cohen's $d = 0.55$, bootstrap CI $[0.45, 0.65]$).

\textbf{Pre-registered prediction:} Essential genes will show higher
variability of dependency profiles across cell lines than non-essential
genes (one-tailed, $d > 0.35$). Framed as confirmatory replication of
the Perturb-seq finding.

\textbf{Result:} $d = 2.65$ (bootstrap CI $[2.52, 2.79]$).
Pre-registered criterion satisfied. PC1 removal changes effect by
0.2\%, ruling out the fitness confound.

\subsection{Discovery of the magnitude artifact}
\label{sec:artifact}

\textbf{Status: EXPLORATORY. Not pre-registered. Discovered post-hoc
after observing that $d = 2.65$ was implausibly large.}

The $d = 2.65$ result triggered scrutiny: Cohen's $d$ values above 2
imply near-zero distributional overlap, which is unusual for biological
comparisons. Two checks revealed the effect is a magnitude artifact.

\subsubsection{DepMap: coefficient-of-variation normalization}

Essential genes have large negative Chronos scores (by definition---this
is what ``essential'' means in DepMap). Larger absolute scores
mechanically produce larger absolute variance. Normalizing by mean
effect magnitude (coefficient of variation: SD / $|\bar{x}|$) removes
this confound:

\begin{itemize}
  \item Essential CV: mean $= 0.38$, median $= 0.26$
  \item Non-essential CV: mean $= 2.95$, median $= 1.71$
  \item Cohen's $d$ (CV): $-1.33$ (sign reversal)
\end{itemize}

After normalization, non-essential genes show \emph{higher} relative
variability. This reversal is biologically interpretable: non-essential
genes are dispensable under most conditions but occasionally recruited
for context-specific functions, producing high relative variability
around a near-zero mean. The raw $d = 2.65$ was the tautology that
``things with bigger effects have bigger variance,'' masking a real (and
opposite-signed) pattern of context-dependent non-essential gene
recruitment.

\subsubsection{Perturb-seq: magnitude-matched comparison}

Geodesic distance correlates with CERES effect magnitude (Spearman $\rho
= 0.33$). After matching essential and non-essential
genes by response magnitude ($n = 127$ matched pairs within 0.2 CERES
units):

\begin{itemize}
  \item Essential mean distance: 2.659
  \item Non-essential mean distance: 2.671
  \item Cohen's $d$: $-0.06$ ($p = 0.62$, Mann-Whitney)
\end{itemize}

The effect vanishes entirely. The original $d = 0.55$ was driven by
essential genes having larger transcriptomic effects, which improves PCA
subspace estimation quality and shifts geodesic distances independent of
true directional diversity.

\subsubsection{Confirming the confounder with a magnitude-invariant metric}

We computed PC1 direction instability---the cosine distance between
unit-normalized primary response directions in K562 and RPE1---for each
gene. Raw PC1 direction instability shows a small essential/non-essential
difference ($d = 0.19$, $p = 0.001$, magnitude correlation $\rho =
0.10$). After magnitude matching ($n = 127$ pairs):

\begin{itemize}
  \item Essential PC1 D (matched): $0.505$
  \item Non-essential PC1 D (matched): $0.512$
  \item Cohen's $d$: $-0.05$ ($p = 0.63$)
\end{itemize}

After residualization (all 1{,}567 genes): $d = 0.07$, $p = 0.19$.

The magnitude-invariant metric converges to the same null as the
magnitude-sensitive one. This rules out the mediator interpretation: if
essential genes had genuinely more diverse response \emph{directions}
(independent of response strength), the cosine-based metric would detect
this regardless of magnitude matching. It does not. The confound
operates through estimation quality at every level.

\subsubsection{Overriding the pre-registered confirmation}

The pre-registered criterion ($d > 0.35$) was technically satisfied by
the raw SD metric ($d = 2.65$). We override this on honesty grounds: the
metric captured magnitude, not the directional divergence it was
intended to measure. A pre-registration protects against cherry-picking
hypotheses; it does not oblige one to report a confounded metric as if
it measures the intended construct. Both the ``confirmation'' and the
artifact are reported; the latter is our actual conclusion.

\subsection{Cross-cell-type transport distances}
\label{sec:ot_confound}

\textbf{Status: EXPLORATORY. Discovered post-hoc after running
pre-registered transport tests on Perturb-seq data.}

Two pre-registered tests asked whether CRISPR knockdown responses for
the \emph{same gene} produce more similar cross-cell-type distances
than knockdowns of \emph{different genes}, testing ``gene-specific
transport'' between K562 and RPE1 \citep{replogle2022mapping}.
The first used geodesic distance on $\text{Gr}(5, 2073)$ between
top-5 PCA subspaces (1{,}676 genes); the second used earth mover's
distance (Wasserstein-1) between full perturbed-cell distributions
(1{,}875 genes, subsampled to 200 cells, 10 repeats, exact POT solver).

Both naive results showed same-gene distances lower than random-gene
distances: geodesic $d = -0.14$ ($p = 2.6 \times 10^{-6}$), EMD
$d = -0.172$ ($p = 9.8 \times 10^{-19}$). Both fell in the
pre-registered ``ambiguous'' range ($|d| < 0.2$).

Both metrics correlate strongly with minimum cell count per gene
(geodesic: Spearman $\rho = -0.84$; EMD: $\rho = -0.47$). Stratifying
by sample size reveals the same bin-inversion pattern as the DepMap
artifact:

\begin{table}[h]
\centering
\caption{Cross-cell-type distance effect size inverts across
sample-size bins for both metrics. Each cell reports Cohen's $d$ of
same-gene distance relative to the random-gene null.}
\label{tab:transport_bins}
\begin{tabular}{lccrr}
\toprule
$\min(n)$ bin & $n_{\text{geo}}$ & $n_{\text{EMD}}$ & Geodesic $d$ & EMD $d$ \\
\midrule
$[20, 40)$   & 204 & 403  & $+1.06$ & $+0.44$ \\
$[40, 60)$   & 407 & 407  & $+0.58$ & $+0.02$ \\
$[60, 100)$  & 575 & 575  & $-0.14$ & $-0.27$ \\
$[100, 200)$ & 429 & 429  & $-1.13$ & $-0.71$ \\
$[200, 700)$ & 61  & 61   & $-2.02$ & $-0.75$ \\
\bottomrule
\end{tabular}
\end{table}

The monotonic inversion is the diagnostic: if gene-specific transport
were biological, all bins would show $d < 0$ regardless of cell count.
Instead, both metrics transition smoothly from positive to negative as
$n$ increases. PCA subspace estimation is more $n$-sensitive than
distributional EMD, so the metric one would naively trust more
(subspace geodesic) is more contaminated.

\section{Causal DAG}
\label{sec:dag}

The magnitude confound has the same causal structure across all
affected metrics, though the mechanism differs.

\subsection{DepMap (variance pathway)}

\begin{equation*}
\text{Magnitude} \rightarrow
\begin{cases}
\text{Essentiality label} & \text{(confounds grouping)} \\
\text{Score variance} & \text{(inflates metric)}
\end{cases}
\end{equation*}

Essentiality is \emph{operationally defined} by dependency-score
magnitude (Chronos $< -0.5$). The magnitude--essentiality association is
near-tautological. Magnitude also causes higher variance (larger
absolute numbers span larger ranges). Adjusting for magnitude is correct
because it removes a spurious path, not a mediating one.

\subsection{Perturb-seq (estimation pathway)}

\begin{equation*}
\text{Magnitude} \rightarrow
\begin{cases}
\text{Essentiality status} & \text{(confounds grouping)} \\
\text{Subspace estimation quality} \rightarrow \text{Geodesic distance}
  & \text{(inflates metric)}
\end{cases}
\end{equation*}

The evidence that this is a confounder rather than a mediator: if
magnitude were a \emph{mediator}, then conditioning on
magnitude would attenuate but not eliminate the essential/non-essential
difference. Instead, magnitude-matching eliminates it completely across
all tested metrics: geodesic distance (matched $d = -0.06$), PC1 cosine
D (matched $d = -0.05$), and PC1 cosine D residualized over all genes
($d = 0.07$, $p = 0.19$).

\subsection{Cross-cell-type transport (estimation pathway)}

\begin{equation*}
\text{Cell count} \rightarrow
\begin{cases}
\text{Same-gene pairing structure} & \text{(confounds grouping)} \\
\text{Estimation quality} \rightarrow \text{Distance metric}
  & \text{(deflates metric)}
\end{cases}
\end{equation*}

Cell count determines estimation quality for both subspace-based
(geodesic, $\rho = -0.84$) and distribution-based (EMD, $\rho = -0.47$)
metrics. Same-gene pairs preserve the within-gene cell count correlation
between K562 and RPE1, while random-gene pairs break this matching.

\subsection{LINCS direction instability (no confound)}

\begin{equation*}
\text{Magnitude} \not\rightarrow D \quad \text{(cosine normalization
blocks the path)}
\end{equation*}

Direction instability divides by $\|\mathbf{s}_i\|$ before computing
angles. The mechanical path from magnitude to metric is blocked.
In finite samples, a residual empirical correlation persists ($\rho
\approx -0.41$ between $D$ and signal $\ell_2$ norm in LINCS): weaker
signals produce noisier unit vectors, slightly inflating $D$. This
second-order effect is documented but small relative to the primary
confound in subspace metrics ($\rho = 0.33$).


\section{Part II: Direction Instability Captures Real Biology}

With magnitude confounds eliminated from contention, we present evidence
that direction instability captures genuine biological signal: it
predicts mechanism transport, stratifies by mechanism of action, and
reveals fine-grained mechanistic gradients.

\subsection{The direction instability landscape}

Across 8{,}949 LINCS drugs tested in $\geq 3$ cell lines, direction
instability spans 0 (identical signatures in every context) to 1.33
(near-orthogonal signatures). The distribution is unimodal with a long
right tail. Drugs with $D < 0.2$ produce consistent transcriptomic
effects regardless of cellular background; those with $D > 0.8$ produce
context-dependent, divergent responses.

\subsection{Gene-level validation: Jaccard overlap confirms directional signal}

For drugs with low $D$ (consistent directions), the top 50 most
perturbed genes should overlap substantially across cell lines. For
drugs with high $D$ (divergent directions), the perturbed gene sets
should differ. Mean top-gene Jaccard similarity correlates with direction
instability at $\rho = -0.72$, confirming that the
geometric metric captures a biological signal: directionally stable drugs
perturb the same genes regardless of context.

\subsection{MOA stratification in LINCS}

Drugs targeting fundamental cellular machinery---HDAC inhibitors,
proteasome inhibitors, mTOR inhibitors, topoisomerase
inhibitors---exhibit low direction instability (mean $D = 0.28$, $n =
1{,}127$). These targets are conserved across cell types, so the
transcriptomic response to their inhibition is consistent. Receptor
agonists, GPCR modulators, and ion channel drugs exhibit high direction
instability (mean $D = 0.71$, $n = 463$). These targets are
differentially expressed across cell types, so the downstream
transcriptomic response depends on which receptors are present.
The separation is large (Cohen's $d = 1.94$, 95\% CI $[1.61, 2.33]$).

\subsection{HDAC inhibitor selectivity gradient}

Among HDAC inhibitors, direction instability reveals a selectivity
gradient invisible to other metrics. Pan-HDAC inhibitors
(vorinostat, panobinostat, belinostat) have the lowest $D$ values ($D
\approx 0.15$--$0.20$): they inhibit all HDAC isoforms and produce
near-identical transcriptomic effects in every cell line. Class-selective
inhibitors (entinostat, tubacin) have intermediate $D$ ($D \approx
0.25$--$0.35$): their effects depend partly on which HDAC isoforms each
cell line expresses. Isoform-selective inhibitors targeting
context-dependent HDACs (PCI-34051 targeting HDAC8
\citep{balasubramanian2008hdac8}) have the highest $D$ values among
HDAC drugs ($D \approx 0.45$).

This gradient is predicted by known biology \citep{bolden2006anticancer}:
pan-inhibitors hit all targets regardless of expression context, while
selective inhibitors depend on context-specific target availability. The
correlation between target breadth and directional consistency confirms
that direction instability captures a mechanistically interpretable
signal.

\subsection{Imatinib case study}

Imatinib \citep{deininger2005imatinib} has low direction instability ($D
= 0.18$, 50th percentile) and high cross-cell-line Jaccard overlap
(0.42, top 5\%). Its primary target, BCR-ABL, drives a stereotyped
transcriptomic response: the same downstream genes are perturbed in
every cell line that expresses the target. Cell lines lacking BCR-ABL
show minimal response (near-zero signatures), which cosine DI handles
correctly---near-zero vectors are treated as uninformative rather than
``different.''

\subsection{Predictive validation: leave-one-out AUROC}
\label{sec:predictive}

If direction instability captures mechanism transport, low-$D$ drugs
should produce consistent signatures even in cell lines excluded from
the original calculation. We test this with leave-one-out
cross-validation: for each drug with $\geq 10$ cell lines ($n =
2{,}694$), we compute $D$ from $K-1$ cell lines, then ask whether the
held-out signature is consistent (cosine similarity $> 0.5$ with the
remaining signatures).

Direction instability predicts the held-out consistency with AUROC $=
0.986$ (95\% CI $[0.983, 0.989]$, $n = 2{,}694$). For each drug, the
held-out cell line is excluded from both the $D$ computation and the
consistency evaluation, so the predictor and outcome share no cell-line
data. A permutation null (shuffling drug labels) yields AUROC $= 0.500
\pm 0.012$ ($n = 1{,}000$ permutations), confirming the signal is
drug-specific. This cosine-matched AUROC is an exploratory result;
the pre-registered construction-neutral benchmark
(Section~\ref{sec:benchmark}) reports AUROC $= 0.945$ $[0.937, 0.952]$
on a Jaccard outcome that cannot benefit from cosine--cosine matching.

\subsection{Robustness controls}

\textbf{Replicate noise:} Computing $D$ from biological replicates
within the same cell line yields near-zero values ($D < 0.05$),
confirming that technical noise does not inflate the metric.

\textbf{Cytotoxicity:} Cytotoxic drugs might produce consistent
signatures (low $D$) by killing cells uniformly, yielding a trivial
``death program'' rather than mechanism-specific transport. Three
controls address this. First, restricting to sub-lethal concentrations
preserves the AUROC (0.979 vs.\ 0.986). Second, removing known
cytotoxicity genes from the 978-landmark set does not change rankings
($\rho > 0.99$, Table~\ref{tab:robustness}). Third, known cytotoxic
drugs span the full $D$ range (mean $D = 0.42$), inconsistent with a
universal low-$D$ death signal.

\textbf{Population null:} Each drug's $D$ was tested against a
population null---the expected $D$ from random drug pairs at the same
sample size. The $p$-value distribution is uniform for most drugs, with
enrichment near zero for low-$D$ drugs and near one for high-$D$ drugs,
matching the expected pattern for a real signal.


\section{Part III: Pre-registered Combined Experiments}

\subsection{Pre-registration integrity chain}

We pre-registered ten decision criteria across three experiments (commit
\texttt{8d74bd0}), specifying quantitative thresholds and interpretation
guides before observing any results. Five criteria passed and five
failed. We report all outcomes against the frozen thresholds, interpret
the failures, and distinguish pre-registered conclusions from post-hoc
observations.

\begin{table}[htbp]
\centering
\caption{Pre-registered criteria outcomes. All thresholds were frozen
before experiments were run (commit \texttt{8d74bd0}). ``Pass'' and
``Fail'' refer strictly to the pre-registered threshold; interpretation
follows in the text.}
\label{tab:prereg_outcomes}
\small
\begin{tabular}{llllr}
\toprule
\textbf{Exp.} & \textbf{Criterion} & \textbf{Threshold} & \textbf{Observed} & \textbf{Result} \\
\midrule
1 & C1: Aggregate $|d|$ & $< 0.1$ & $d = -0.22$ & Fail \\
1 & C2: Bin-flatness & All bins $\in [-0.15, +0.15]$ & $-0.19$ to $-0.23$ & Fail \\
1 & C3: Cell-count $|\rho|$ & $< 0.20$ & $\rho = -0.16$ & Pass \\
\midrule
2 & C1: DI cosine gap & $\geq 0.02$ over alternatives & $-0.007$ & Fail \\
2 & C2: DI Jaccard gap & $> 0$ over alternatives & $+0.088$ & Pass \\
2 & C3: Gap shrinks & Cosine gap $>$ Jaccard gap & Reversed & Fail \\
2 & C4: Correction converges & Corrected gap $<$ raw gap & $0.088 \to 0.023$ & Pass \\
\midrule
3 & C1: Group $d > 0.3$ & Machinery $<$ receptor & $d = 0.77$ & Pass \\
3 & C2: Three-tier ordering & Mach.\ $<$ kinase $<$ receptor & Kinase $<$ mach. & Fail \\
3 & C3: Significance & $p < 0.05$ & $p = 0.002$ & Pass \\
\bottomrule
\end{tabular}
\end{table}

\subsection{Experiment 1: cosine direction instability on Perturb-seq}

This experiment tests the central claim of the magnitude confound
analysis: cosine-based direction instability should show flat bins across
cell-count strata, unlike geodesic distance and EMD, which invert.

Direction instability detected a small, consistent same-gene transport
signal ($d = -0.22$, 95\% CI $[-0.33, -0.19]$; same-gene mean $D =
0.152$, random-gene mean $D = 0.161$). Both C1 and C2 failed their pre-registered thresholds: the
aggregate $|d|$ exceeded 0.1, and per-bin effect sizes fell outside the
$\pm 0.15$ band. The pre-registered scientific hypothesis, however,
was immunity to bin inversion---the pattern that invalidates confounded
metrics. Across five cell-count bins spanning $[20, 700)$ cells,
per-bin effect sizes ranged from $-0.19$ to $-0.23$: a total span of
0.04. For comparison, geodesic distance spans 3.08 across the same bins
($d = +1.06$ to $d = -2.02$); earth mover's distance spans 1.19
($+0.44$ to $-0.75$). The cosine metric shows a consistent shift
rather than inversion, and the cell-count correlation is small ($\rho =
-0.16$, C3 passed). The $\pm 0.15$ auxiliary band did not anticipate a
genuine weak same-gene signal at this magnitude, but the absence of
inversion---the primary diagnostic---is confirmed.
We report the absolute means and effect size rather than the $p$-value
($2.8 \times 10^{-13}$), since a significant-but-small effect at $n =
1{,}676$ is precisely the pattern this paper warns against
overinterpreting.

\subsection{Experiment 2: multi-metric LOO benchmark on LINCS}
\label{sec:benchmark}

This experiment evaluates direction instability against three alternative
metrics (Fr\'{e}chet variance, magnitude CV, mean top-gene Jaccard) on
the predictive benchmark from Section~\ref{sec:predictive}, using a
construction-neutral outcome (Jaccard overlap) to avoid the structural
advantage that cosine--cosine matching would confer.

Direction instability achieved the highest AUROC on the
construction-neutral Jaccard outcome ($0.945$, 95\% CI $[0.937, 0.952]$
vs.\ $0.857$ $[0.841, 0.872]$ for Fr\'{e}chet variance; paired
$\Delta = +0.088$ $[+0.075, +0.102]$, $p < 10^{-4}$; C2 passed). After
magnitude residualization, DI retained its advantage (corrected AUROC
$0.754$ $[0.735, 0.773]$ vs.\ $0.731$ $[0.712, 0.751]$; C4 passed). On
the cosine outcome, Fr\'{e}chet variance ($0.976$ $[0.968, 0.983]$)
exceeded DI ($0.969$ $[0.958, 0.979]$; paired $\Delta = -0.007$
$[-0.012, -0.002]$, $p = 0.002$; C1 failed).

C3 failed because DI's advantage over alternatives was \emph{larger} on
the Jaccard outcome than on cosine---the opposite of our prediction that
the construction-matched advantage would shrink. This reversal is
consistent with the paper's central thesis: the cosine outcome is
defined by the same geometry the matched metric optimizes, so
construction matching flatters the matched metric. DI's advantage
emerges on the construction-neutral outcome, where it cannot benefit
from geometric self-similarity. The two failures strengthen rather than
weaken the central claim, but we report them as failures of the
pre-registered predictions.

\subsection{Experiment 3: MOA class stratification in JUMP Cell Painting}

This experiment tests whether the LINCS MOA stratification pattern
(Section~5.3) replicates in a different measurement modality: Cell
Painting morphological profiles.

Drugs targeting fundamental cellular machinery showed lower direction
instability in Cell Painting morphological profiles (mean $D = 0.69$,
$n = 18$) than receptor-mediated drugs (mean $D = 0.83$, $n = 88$;
$d = 0.77$, 95\% CI $[0.33, 1.35]$, Mann--Whitney $p = 0.002$; C1 and C3 passed). This
replicates the LINCS L1000 pattern (LINCS $d = 1.94$) in a different
measurement modality, with expected attenuation.

The three-tier ordering failed (C2): kinase inhibitors (mean $D = 0.65$,
$n = 21$) fell slightly below machinery rather than between machinery
and receptor. At $n = 18$ and $n = 21$, the $0.04$ mean difference
between machinery and kinase is well within sampling noise; the
distinction between these two groups is underpowered. The robust result
is the machinery-vs-receptor separation, which holds with a large
effect size across measurement technologies.

\section{The Fix}

Three normalizations remove the magnitude confound from affected
metrics:

\begin{enumerate}
  \item \textbf{Coefficient of variation} (SD / $|\text{mean effect}|$):
    removes scale dependence. Applicable when effects have non-zero means.
  \item \textbf{Magnitude matching}: compare only perturbations with
    similar absolute effect sizes. The most assumption-free approach; does
    not require linearity.
  \item \textbf{Residualization}: regress out effect magnitude and test
    residuals. Assumes a linear (or specified) relationship between
    magnitude and the metric.
\end{enumerate}

Cosine-based formulations (direction instability, Equation~1) require no
correction---they are the magnitude-normalized version by construction.


\section{Discussion}

This paper makes two claims. First, variance-based, subspace-based, and
distribution-based perturbation-divergence metrics are
magnitude-confounded when used to compare groups that differ in effect
strength or sample size. Second, cosine-based direction instability is
immune to this confound and captures genuine biological signal: mechanism
transport, MOA stratification, and mechanistic selectivity gradients.

The magnitude confound claim is supported by three instances of the same
discovery pattern: a pre-registered metric yields an apparently
meaningful effect (DepMap $d = 2.65$, Perturb-seq geodesic $d = 0.55$,
cross-cell-type transport $d = -0.14$ to $-0.172$), which post-hoc
analysis reveals as estimation-quality confounding. In each case, the
signal inverts or vanishes after controlling for the confound.

The direction instability claim is supported by predictive validation
(LOO AUROC $= 0.945$ $[0.937, 0.952]$ on the construction-neutral
outcome; $0.986$ $[0.983, 0.989]$ on the exploratory cosine-matched
outcome), cross-technology MOA replication (LINCS $d = 1.94$, JUMP-CP
$d = 0.77$), and pre-registered experiments where five of ten criteria
passed. The criteria that failed are interpretable: the Perturb-seq
bin-flatness band did not anticipate a genuine weak transport signal
(bins span 0.04, vs.\ 3.08 for geodesic---no inversion, but shifted
outside the $\pm 0.15$ band), Fr\'{e}chet variance exceeded DI on the
cosine outcome by 0.007 AUROC (paired $p = 0.002$), and the three-tier
MOA ordering was underpowered at $n = 18$--$21$.

\subsection{Pre-registration as diagnostic}

Half of our pre-registered thresholds missed. We treat this as
informative rather than embarrassing. The failures reveal where our
assumptions were wrong: the Perturb-seq bin-flatness threshold assumed
the cosine metric would show near-zero aggregate effect, when in fact a
genuine weak signal exists; the cosine--Jaccard gap prediction assumed
DI's advantage would be construction-dependent, when it turned out to
be construction-independent; Fr\'{e}chet variance's marginal win on the
cosine outcome ($\Delta = 0.007$) is consistent with the
construction-matching critique the paper advances. Each failed
prediction is itself a finding.

The practice of overriding a technically-passed criterion (DepMap $d =
2.65$) on honesty grounds deserves emphasis. A pre-registration
protects against cherry-picking hypotheses; it does not oblige one to
report a confounded metric as if it measures the intended construct.
Rigid adherence to pre-registered decisions when the decision criteria
themselves are revealed as inadequate serves pre-registration theater,
not scientific integrity.

\subsection{Implications for perturbation biology}

The practical implication is a rule: before comparing ``divergence''
across groups that differ in perturbation strength or cell count,
normalize for magnitude. Failing to do so produces effects that look
large ($d = 2.65$, $d = 0.55$) or directionally misleading ($d = -0.14$
averaging over bins from $+1.06$ to $-1.24$) but reduce to the
tautology that bigger effects are noisier or better-estimated.

Researchers computing variability metrics on perturbation atlases
(DepMap, Perturb-seq, PRISM, JUMP-CP) should verify that their metric is
magnitude-invariant before interpreting group differences as biological
divergence. The transport confound is particularly insidious because it
can masquerade as a weak positive signal: the aggregate $d = -0.14$
appeared to be a real-but-small biological effect, when it was an
average over bins pointing in opposite directions.

Direction instability offers a constructive alternative. Its magnitude
immunity is mechanical (cosine normalization blocks the path from
magnitude to metric), and its biological validity is confirmed by
predictive performance (construction-neutral LOO AUROC $= 0.945$
$[0.937, 0.952]$), cross-technology replication, and mechanistically
interpretable gradients within drug classes.

\subsection{Limitations}

The magnitude-matching analysis used $n = 127$ matched pairs,
constrained by the overlap in effect-size distributions between
essential and non-essential genes. The matched sample represents the
narrow region where both distributions overlap: atypically weak essential
genes matched to atypically strong non-essential genes. However, the
residualization analysis uses all 1{,}567 labeled genes and converges to
the same null ($d = 0.07$, $p = 0.19$ for PC1 D), confirming that the
matching result is not an artifact of sample restriction.

Our claim that direction instability is ``immune'' to magnitude
confounding is precise: the cosine normalization removes the
\emph{mechanical} dependency. A residual empirical correlation persists
in finite samples ($\rho = 0.10$ for PC1 D vs.\ magnitude in
Perturb-seq; $\rho \approx -0.41$ between $D$ and signal norm in LINCS
drug data). This arises because weaker signals produce noisier estimated
directions, slightly inflating $D$. The metric is mechanically immune but
not entirely empirically independent of magnitude when directions are
estimated from finite data.

The JUMP Cell Painting MOA stratification ($n = 18$ machinery, $n = 21$
kinase, $n = 88$ receptor) is underpowered for fine-grained ordering.
The machinery-vs-receptor separation is robust ($d = 0.77$, $[0.33,
1.35]$, $p = 0.002$), but the kinase-machinery distinction ($\Delta =
0.04$) is not
resolvable at these sample sizes.

The Perturb-seq analysis uses only two cell types (K562, RPE1), so
direction instability reduces to a single pairwise cosine rather than
a population average. This suffices for the bin-flatness test (the
diagnostic question is whether per-bin $d$ values invert with cell
count, not whether the absolute level of transport is high) but limits
generalization to multi-context settings.

\section{Conclusion}

We pre-registered divergence and transport predictions, naive metrics
confirmed them, and our own post-hoc checks killed them---three times
with the same pattern. The killing is the finding: geometric
perturbation metrics measure magnitude and sample size unless explicitly
prevented from doing so. The DAG is simple, the fix is simple, and the
failure mode is likely widespread in any analysis comparing perturbation
groups that differ in effect strength or cell count.

The constructive counterpart: direction instability, which blocks the
magnitude path by construction, predicts cross-cell-line mechanism
transport (construction-neutral LOO AUROC $= 0.945$ $[0.937, 0.952]$),
replicates MOA stratification across measurement technologies, and
resolves mechanistic gradients within drug classes. The metric works
because it asks the right question---do perturbation effects point in the
same direction?---rather than a question contaminated by how precisely
we estimated them.

\section*{Data and Code Availability}

All analysis code and pre-registration documents are available at
\url{https://github.com/elliottower/bracket-norm-validity} and
\url{https://github.com/elliottower/drug-perturbation-geometry}.
The repositories contain the full commit history demonstrating the
pre-registration integrity chains (commits \texttt{1dc20a2},
\texttt{249abaf}, \texttt{98897ed}, \texttt{92e3276}, \texttt{8d74bd0}).
Zenodo DOIs: drug transport dataset
(\texttt{zenodo.org/records/21207728}), magnitude confound replication
materials (\texttt{zenodo.org/records/21207429}).
DepMap data (24Q4 release) is available from
\url{https://depmap.org/portal/}. LINCS L1000 data is from GEO
(GSE92742). Perturb-seq data is from \citet{replogle2022mapping}. JUMP
Cell Painting data is from \citet{chandrasekaran2023jump}.

\section*{Author Contributions}

E.T.\ conceived the study, designed experiments, wrote
pre-registrations, implemented all analyses, discovered the magnitude
artifact, and wrote the manuscript.

\section*{AI Disclosure}

Claude Code (Anthropic) was used as a coding assistant for data
analysis, statistical computation, and manuscript preparation.
Perplexity Pro was used for literature search and as an independent
reviewer of pre-registration documents and experimental scripts.
All scientific decisions, experimental design, interpretations, and
conclusions are the author's.

\bibliographystyle{plainnat}
\bibliography{references}

\appendix

\section{MOA Classification Statistics}
\label{app:moa}

\begin{table}[h]
\centering
\caption{MOA class sizes used in stratification analyses. LINCS
classification uses Drug Repurposing Hub annotations. JUMP-CP
classification uses the same MOA lists applied to InChIKey-matched
compounds.}
\begin{tabular}{lrr}
\toprule
MOA class & LINCS $n$ & JUMP-CP $n$ \\
\midrule
Machinery (HDAC, proteasome, mTOR, topoisomerase, \ldots) & 1{,}127 & 18 \\
Kinase inhibitors & 892 & 21 \\
Receptor (GPCR, ion channel, nuclear receptor, \ldots) & 463 & 88 \\
Other / unclassified & 6{,}467 & 203 \\
\bottomrule
\end{tabular}
\end{table}

\section{HDAC Inhibitor Selectivity Gradient}
\label{app:hdac}

\begin{table}[h]
\centering
\caption{Direction instability of HDAC inhibitors in LINCS L1000, ranked
by $D$. Pan-HDAC inhibitors cluster at low $D$; isoform-selective
inhibitors at high $D$.}
\begin{tabular}{llccc}
\toprule
Drug & Selectivity & $D$ & $n_{\text{cells}}$ & Clinical phase \\
\midrule
Panobinostat & Pan-HDAC & 0.14 & 64 & Approved \\
Vorinostat & Pan-HDAC & 0.16 & 63 & Approved \\
Belinostat & Pan-HDAC & 0.19 & 42 & Approved \\
Entinostat & Class I & 0.27 & 38 & Phase III \\
Tubacin & HDAC6 & 0.32 & 25 & Tool compound \\
PCI-34051 & HDAC8 & 0.45 & 18 & Preclinical \\
\bottomrule
\end{tabular}
\end{table}

\section{Top 50 Directionally Stable Drugs}
\label{app:top50}

The 50 drugs with lowest direction instability among those tested in
$\geq 10$ cell lines are dominated by: pan-HDAC inhibitors,
proteasome inhibitors (bortezomib, carfilzomib), HSP90 inhibitors
(tanespimycin, alvespimycin), and topoisomerase inhibitors (irinotecan,
topotecan). All are cytotoxic agents targeting highly conserved cellular
machinery, consistent with the mechanism-transport interpretation.
The full ranked list with $D$ values, $n$ cell lines, MOA, and targets
is provided in the Zenodo supplement.

\end{document}
